\section{Decomposition theorem in the Poisson geometric setting}

\label{section:geometric-setting}

\subsection{The coadjoint module for a Lie algebra}

Let $d \geqslant 2$ be an integer. In this paper, $\g$ will denote the Lie algebra $\sl_d$ of square matrices of size $d$ whose trace equals $0$. Denote by $x_{i,j}$ the elementary matrix whose coefficients are all $0$, except in position $(i,j)$ where it is one. We identify the dual vector space $\g^\vee$ with the affine Poisson variety
\[ \g^\vee \coloneqq \Spec (\Sym \g), \]
where the Poisson bracket on $\Sym \g$ extends the Lie bracket on $\g$.

Denote by $G$ the adjoint group of $\g$: it is the connected component of the identity in the group of Lie algebra automorphisms of $\g$. Explicitly, it is isomorphic to the projective linear group $\mathrm{PGL}_d$. Denote by $\cdot$ the adjoint action of $G$ on $\g$ and the coadjoint action on $\g^\vee$. The action on $\g^\vee$ is by Poisson automorphisms.


Denote by 
\[ (x \mid y) \coloneq \operatorname{trace}(xy), \quad \text{for} \quad x, y \in \g,  \]
the normalised Killing form on $\g$. It is $G$-invariant and non-degenerate, therefore it induces a $G$-equivariant isomorphism
\[ \kappa \colon \g \overset{\sim}{\longrightarrow} \g^\vee, \quad x \longmapsto (x \mid \bullet). \]

\subsection{The Hamiltonian action associated with the minimal nilpotent orbit}

We recall the classical construction of the Poisson structure on the minimal Slodowy slice. It is a particular case of a more general construction for any nilpotent orbit \cite{kostant1978whittaker,lynch1979generalized,premet2002special,gan2002quantization,losev2010quantized}. 

Consider the minimal $\sl_2$-triple of $\g$ given by
\[ e \coloneqq x_{1,2}, \quad h \coloneqq x_{1,1} - x_{2,2} \quad \text{and} \quad f \coloneq x_{2,1}. \]
We consider the diagonal element $H \coloneqq x_{1,1} - \frac{1}{d} \sum_{i=1}^{d} x_{i,i}$. Because it is semisimple, its adjoint action defines a Lie algebra grading on $\g$. It is explicitly given by:
\[ \g_{-1} = \bigoplus_{i=2}^d \C x_{i,1}, \quad \g_0 = \bigoplus_{i,j=2}^d \C (x_{i,j} - \delta_{i=j}x_{1,1}), \quad \g_{1}  = \bigoplus_{j=2}^d \C x_{1,j}. \]

It is an even good $\frac{1}{2} \Z$-grading in the sense of \cite{elashvili2005classification,brundan2007good}. In particular, the adjoint action of $f$ gives an injective map $\g_1 \hookrightarrow \g_0$ and a surjective map $\g_0 \twoheadrightarrow \g_{-1}$.

We denote the dual grading by
\[ \g^\vee = (\g^\vee)_1 \oplus (\g^\vee)_0 \oplus (\g^\vee)_{-1}, \]
where the subspace $(\g^\vee)_i \subseteq \g^\vee$ is defined as the image of the dual of the projection $\g \twoheadrightarrow \g_{-i}$ for all $i$ in $\Z$. In particular $(\g^\vee)_i \cong (\g_{-i})^\vee$.

The positive part of the grading $\g_{\geqslant 1} = \g_1$ is an Abelian nilpotent subalgebra of $\g$. It is the Lie algebra of the following unipotent subgroup of $G$:
\[ G_{\geqslant 1} \coloneqq \bigg\{\left(\begin{smallmatrix}
	1 & \star & \cdots & \star \\
	& 1 & & (0) \\
	& & \ddots & \\
	&(0)  & & 1
\end{smallmatrix}\right)\bigg\}. \]


Denote by $\chi \coloneq \kappa(f)$ the linear form on $\g$ associated to $f$ and by $\overline{\chi} \coloneqq \chi|_{\g_{\geqslant 1}}$ its restriction to $\g_{\geqslant 1}$. Because this Lie algebra is Abelian, $\overline{\chi}$ is a character. The restriction map $\g^\vee \twoheadrightarrow (\g_{\geqslant 1})^\vee$ is a moment map, and we can twist it by the character $\overline{\chi}$, leading to the moment map
\[ \mu \colon \g^\vee \longrightarrow (\g_{\geqslant 1})^\vee, \quad \xi \longmapsto \xi|_{\g_{\geqslant 1}} + \overline{\chi}. \]
Its fibre over $0$ is the following affine subspace of $\g^\vee$:
\[ \mu^{-1}(0) = - \chi + (\g^\vee)_{\geqslant 0}. \]

To compute its quotient modulo the action of $G_{\geqslant 1}$, one introduces the Slodowy slice associated to the minimal $\sl_2$-triple:
\[ S_f \coloneqq - \chi + [\g, e]^\perp = \kappa(-f + \g^e). \]
The Slodowy slice is a closed subvariety of $ \mu^{-1}(0)$ because $\g^e$ is non-negatively graded.

\begin{theorem}[{\cite[Lemma 2.1]{gan2002quantization}}] \label{theorem:gan-ginzburg-isomorphism}
	The action map
	\[ \phi \colon G_{\geqslant 1} \times S_f \longrightarrow - \chi + (\g^\vee)_{\geqslant 0}, \quad (g, \xi) \longmapsto g \cdot \xi  \]
	is a $G_{\geqslant 1}$-equivariant isomorphism of algebraic varieties. By taking quotient on both sides with respect to the $G_{\geqslant 1}$-actions, one gets an isomorphism
	\[ S_f \cong \mu^{-1}(0) \git G_{\geqslant 1}. \]
\end{theorem}

\subsection{Restriction to a distinguished open set}

Consider the distinguished open subscheme
\[ \g^\vee[e^{-1}] \coloneqq \Spec ((\Sym \g)[e^{-1}]). \]
It is a Poisson variety, stable by the action of $G_{\geqslant 1}$ because $e$ is $G_{\geqslant 1}$-invariant, and we have the induced moment map given by restriction of $\mu$:
\[ \mu_{e^{-1}} \colon \g^\vee[e^{-1}] \longrightarrow (\g_{\geqslant_1})^\vee. \]

\begin{proposition} \label{proposition:hamiltonian-reduction-open}
	The natural open inclusion 
	\[ (\mu_{e^{-1}})^{-1}(0) \hooklongrightarrow \mu^{-1}(0) \]
	descends to a Poisson isomorphism between the corresponding Hamiltonian reductions:
	\[ (\mu_{e^{-1}})^{-1}(0) \git G_{\geqslant 1} \cong \mu^{-1}(0) \git G_{\geqslant 1} \]
	
\end{proposition}

Concretely, this proposition relies on the easy observation that $\mu^{-1}(0)$ is contained in the distinguished open subscheme $\g^\vee[e^{-1}]$. Because later on we will deal with localisation in noncommutative settings, we give details for the proof in term of commutative algebra.

\begin{proof}
	Denote by $I_{\mu}$ the ideal of $\Sym \g$ that is generated by the elements $x_{1,2} + 1$ and $x_{1,i}$ for $3 \leqslant i \leqslant d$. Then:
	\[ (\mu_{e^{-1}})^{-1}(0) = \Spec ((\Sym\g) [e^{-1}] / I_{\mu} [e^{-1}]) \quad \text{and} \quad \mu^{-1}(0) = \Spec (\Sym\g / I_{\mu}). \]
	Because localisation is exact \cite[Proposition 2.5]{eisenbud1995commutative}, one has a natural algebra isomorphism
	\[ (\Sym\g) [e^{-1}] / I_{\mu} [e^{-1}] \cong ((\Sym\g) / I_{\mu}) [e^{-1}] \]
	
	
	The image of $e = x_{1,2}$ in the quotient $(\Sym\g) / I_{\mu}$ is equal to $-1$, so it is invertible. Therefore, the localisation homomorphism
	\[ (\Sym\g) / I_{\mu} \longrightarrow ((\Sym\g) / I_{\mu}) [e^{-1}] \]
	is an isomorphism by \cite[Exercise 2.9]{eisenbud1995commutative}. This shows the isomorphism
	\[ (\mu_{e^{-1}})^{-1}(0) \cong \mu^{-1}(0). \]
	
	This isomorphism is clearly $G_{\geqslant 1}$-equivariant, so we get isomorphic quotients. This isomorphism is Poisson because it is induced by the $G_{\geqslant 1}$-equivariant Poisson embedding $\g^\vee[e^{-1}] \hookrightarrow \g^\vee$ that is compatible with the moment maps $\mu_{e^{-1}}$ and $\mu$.
\end{proof}


\subsection{Slodowy slice as the quotient of a distinguished open set} \label{subsection:slodowy-quotient-open}

The Lie algebra $\Z$-grading
\[ \g = \g_{-1} \oplus \g_0 \oplus \g_1 \]
is equivalent to an action of the multiplicative group $\Gm$ on $\g$ by Lie algebra automorphisms. Concretely, we have an algebraic group homomorphism
\[ \tau \colon \Gm \longrightarrow \operatorname{PGL}_d \cong G, \quad t \longmapsto \left(\begin{smallmatrix}
	t & & & 0 \\
	& 1 & & \\
	& & \ddots & \\
	0 & & & 1
\end{smallmatrix}\right) \in \mathrm{PGL}_d, \]
such that for any $i$ in $\{-1, 0, 1\}$, $x$ in $\g_i$ and $t$ in $\Gm$, one has:
\[ \tau(t) \cdot x = t^i x. \]

Define the following Abelian group of matrices:
\[ K \coloneqq \Gm \times \Ga{}^{d-2}. \]
We have the following closed embedding of algebraic group:
\[ \alpha \colon K \longrightarrow \mathrm{PGL}_d \cong G, \quad (t, a_1, \dots, a_{d-2}) \longmapsto \left(\begin{smallmatrix}
	t & & & & (0) \\
	0 & 1 & &  & \\
	0 & a_1 & 1  & & \\
	\vdots & \vdots & & \ddots &  \\
	0 & a_{d-2} & (0) & & 1
\end{smallmatrix}\right). \]
The corresponding Lie algebra $\kfk$ is spanned in $\g$ by $H$ and the $x_{i,2}$ for $3 \leqslant i \leqslant d$.

\begin{lemma} \label{lemma:afk}
	There is a vector space decomposition
	\[ \g_0 = \kfk \oplus (\g_0)^f. \]	
\end{lemma}

\begin{proof}
	A straightforward computation on the generators of $\kfk$ shows that $\kfk \cap (\g_0)^f = 0$. Because the grading is good in the sense of \cite{elashvili2005classification,brundan2007good}, there is an exact sequence
	\[ 0 \longrightarrow (\g_0)^f \longrightarrow \g_0 \overset{[f, \bullet]}{\longrightarrow} \g_{-1} \longrightarrow 0, \]
	therefore $\dim \g_0 = \dim(\g_0)^f + \dim \g_1$. Since $\dim \g_1 = \dim \kfk = d - 1$, we can conclude that the inclusion $\kfk \oplus (\g_0)^f \subseteq \g_0$ is an equality.
\end{proof}

For $1 \leqslant i \leqslant d - 2$ an integer, introduce the linear form $\chi_i \coloneqq \kappa(x_{i+2, 1})$. Then we have the following explicit description for the distinguished open set:
\[ \g^\vee[e^{-1}] = (\C - \{0\}) \chi + \sum_{i=1}^{d-2} \C \chi_i + (\g^\vee)_{\geqslant 0}. \]

\begin{lemma} \label{lemma:dilation-isomorphism}
	The action map
	\[ \begin{split}	
		\phi_0 \colon K \times (- \chi + (\g^\vee)_{\geqslant 0}) & \longrightarrow  \g^\vee[e^{-1}], \\
		((t, a_1, \dots, a_{d-2}), \xi) & \longmapsto  \alpha(t, a_1, \dots, a_{d-2}) \cdot \xi. 
	\end{split} \]
	is an isomorphism of varieties.
\end{lemma}

\begin{proof}
	The action of $(t, a_1, \dots, a_{d-2})$ in $K$ on $f$ is given by the formula
	\[ \alpha(t, a_1, \dots, a_{d-2}) \cdot (-f) = \left(\begin{smallmatrix}
		0 & & & & (0) \\
		-1/t & 0 & &  & \\
		-a_1 / t & 0 & 0  & & \\
		\vdots & \vdots & & \ddots &  \\
		-a_{d-2} / t & 0 & (0) & & 0
	\end{smallmatrix}\right). \]
	By applying $\kappa$, we get an isomorphism
	\[ K \overset{\sim}{\longrightarrow} (\g_{\geqslant 1})^\vee[e^{-1}], \quad (t, a_1, \dots, a_{d-2}) \longmapsto - \alpha(t, a_1, \dots, a_{d-2}) \cdot \overline{\chi}. \]
	Because the projection $\g^\vee \twoheadrightarrow (\g_{\geqslant 1})^\vee$ is $K$-equivariant, the pull-back of this isomorphism is the isomorphism stated in the lemma.
\end{proof}

The previous lemma implies that the adjoint action of the group $K \ltimes G_{\geqslant 1}$ stabilises the open subset $\g^\vee[e^{-1}]$. It acts by Poisson automorphisms, so the categorical quotient
\[ \g^\vee[e^{-1}] \git (K \ltimes G_{\geqslant 1}) \coloneq \Spec \big((\Sym \g)[e^{-1}]\big)^{K \ltimes G_{\geqslant 1}} \] 
inherits a Poisson structure from $\g^\vee$.

The following statement is an analogue of the Gan--Ginzburg isomorphism (\Cref{theorem:gan-ginzburg-isomorphism}), but the codomain is the open subscheme $\g^\vee[e^{-1}]$ instead of the moment map fibre $\mu^{-1}(0)$.

\begin{proposition} \label{proposition:dilation-gan-ginzburg}
	There is a $K \ltimes G_{\geqslant 1}$-equivariant isomorphism of varieties given by:
	\[ \begin{split}
		\phi' \colon (K \ltimes G_{\geqslant 1}) \times S_f & \longrightarrow \g^\vee[e^{-1}], \\
		((t, a_1, \dots, a_{d-2}), g, \xi) & \longmapsto \alpha(t, a_1, \dots, a_{d-2}) g \cdot \xi. 
	\end{split} \]
\end{proposition}

\begin{proof}
	Recall the isomorphism given by \Cref{theorem:gan-ginzburg-isomorphism}:
	\[ \phi \colon G_{\geqslant 1} \times S_f \overset{\sim}{\longrightarrow} - \chi + (\g^\vee)_{\geqslant 0}, \quad (g, \xi) \longmapsto g \cdot \xi.  \]
	Moreover, we have the isomorphism $\phi_0$ given by \Cref{lemma:dilation-isomorphism}. Therefore, we have the explicit formula for an inverse of $\phi'$:
	\[ \phi'^{-1}(\nu) = ((t, a_1, \dots, a_{d-2}), g, \xi) \quad \text{where} \quad \begin{cases}
		(t, a_1, \dots, a_{d-2}, \nu_1) \coloneq \phi_0{}^{-1}(\nu), \\
		(g, \xi) \coloneq \phi^{-1}(\nu_1).
	\end{cases} \]
\end{proof}

We can deduce from the previous proposition a description of the Slodowy slice as the categorical quotient of the distinguished open set.

\begin{theorem} \label{theorem:slodowy-git}
	There is an isomorphism
	\[ \mu^{-1}(0) \git G_{\geqslant 1} \cong \g^\vee[e^{-1}] \git (K \ltimes G_{\geqslant 1}) \]
	of Poisson varieties.
\end{theorem}

\begin{proof}
	Consider the inclusions
	\[ S_f \hooklongrightarrow \mu^{-1}(0) \hooklongrightarrow \g^\vee[e^{-1}] . \]
	After taking quotients, we get homomorphisms
	\[ S_f \longrightarrow \mu^{-1}(0) \git G_{\geqslant 1} \longrightarrow \g^\vee[e^{-1}]\git (K \ltimes G_{\geqslant 1}) \]
	of schemes that are isomorphisms because $\phi$ and $\phi'$ are isomorphisms. 
	
	To check they are Poisson, let us examine their comorphisms. By \Cref{proposition:hamiltonian-reduction-open}, there is a Poisson isomorphism
	\[ (\Sym \g / I_{\mu})^{G_{\geqslant 1}} \cong ((\Sym \g)[e^{-1}] / I_{\mu} [e^{-1}])^{G_{\geqslant 1}}   \]
	induced by localisation.
	
	Moreover, the natural algebra homomorphism
	\[ (\Sym \g)[e^{-1}] \longrightarrow  (\Sym \g)[e^{-1}] / I_{\mu} [e^{-1}] \]
	induces an algebra isomorphism
	\[ \big((\Sym \g)[e^{-1}]\big)^{K \ltimes G_{\geqslant 1}} \cong \big((\Sym \g)[e^{-1}] / I_{\mu} [e^{-1}]\big)^{G_{\geqslant 1}} \]
	after descending to invariant subalgebras. This isomorphism is Poisson because both Poisson brackets are induced by the one on $(\Sym \g)[e^{-1}]$.
\end{proof}

%\begin{corollary}
%	There is an isomorphism
%	\[ \Jet_\infty \mu^{-1}(0) \git \Jet_\infty G_{\geqslant 1} \cong \Jet_\infty \g^\vee[e^{-1}] \git \Jet_\infty(K \ltimes G_{\geqslant 1}) \]
%	of Poisson vertex schemes.
%\end{corollary}
%
%\cite[Theorem 5.1]{bozec2025functorial}
%
%Consider the nilpotent Lie subalgebra
%

\subsection{Decomposition theorem for the minimal Slodowy slice}

Consider the Poisson schemes
\begin{multline*}
	(\kfk \ltimes \g_{\geqslant 1})^\vee[e^{-1}] = \Spec (\Sym(\kfk \ltimes \g_{\geqslant 1})[e^{-1}]) \\ \text{and} \quad S_f \cong \g^\vee[e^{-1}] \git (K \ltimes G_{\geqslant 1}) = \Spec\big((\Sym \g)[e^{-1}]\big)^{K \ltimes G_{\geqslant 1}}. 
\end{multline*}  
Denote by $\overline{\chi}'$ the restriction of the linear form $\chi = (f \mid \bullet)$ to the Lie subalgebra $\kfk \ltimes \g_{\geqslant 1}$.

\begin{lemma} \label{lemma:open-orbit}
	The orbit map of $-\overline{\chi}'$
	\[ K \ltimes G_{\geqslant 1} \longrightarrow (\kfk \ltimes \g_{\geqslant 1})^\vee, \quad g \longmapsto - g \cdot \overline{\chi}' \]
	induces a scheme isomorphism
	\[ K \ltimes G_{\geqslant 1} \cong (\kfk \ltimes \g_{\geqslant 1})^\vee[e^{-1}]. \]
	Therefore $(\kfk \ltimes \g_{\geqslant 1})^\vee[e^{-1}]$ is an open coadjoint orbit, in particular a symplectic smooth variety.
\end{lemma}

\begin{proof}
	We compute the orbit map in stages, starting with the action of $G_{\geqslant 1}$. Because the group is unipotent, we can use the exponential of an element $x$ in $\g_{\geqslant 1}$. If we iterate the coadjoint action of $x$ on $\overline{\chi}'$, then $x \cdot (x \cdot \overline{\chi}') = 0$ because $[x, [x,f]]$ belongs to $\g_1$ and is orthogonal to $\kfk \ltimes \g_{\geqslant 1}$. 
	
	Moreover, because $(\bullet \mid \bullet)$ gives a perfect pairing between $\g_{-1}$ and $\g_1$, and by using the decomposition of \Cref{lemma:afk}, we deduce that 
	\[ (\bullet \mid \bullet) \colon \g_1 \otimes \g_0 \longrightarrow \C, \quad x \otimes y \longmapsto ([x, f] \mid y) = (x \mid [f, y]) \]
	induces a perfect pairing between $[\g_{\geqslant 1}, f]$ and $\kfk$. 
	
	Therefore, the $\g_1 \cong G_{\geqslant 1}$-orbit map induces an isomorphism:
	\[ \g_1 \overset{\sim}{\longrightarrow} -\overline{\chi}' + \kfk^\vee \subseteq (\kfk \ltimes \g_{\geqslant 1})^\vee, \quad x \longmapsto -\overline{\chi}' - x \cdot \overline{\chi}',  \]
	where $\kfk^\vee$ is identified to the subspace $(\g_{\geqslant 1})^\perp$ of $(\kfk \ltimes \g_{\geqslant 1})^\vee$.
	
	If we now look at the action of $K$, we can deduce the lemma from the computation of \Cref{lemma:dilation-isomorphism}.
%	We use the isomorphism $\kappa\colon \g \overset{\sim}{\rightarrow} \g^\vee$ to make the computation.
%	
%	The action of $(t, a_1, \dots, a_{n-2})$ in $A$ on $f$ is explicitly described by the formula
%	\[ \left(\begin{smallmatrix}
%		t & b_1 & b_2& \cdots & b_{d-1} \\
%		0 & 1 & &  & \\
%		0 & a_1 & 0  & & \\
%		\vdots & \vdots & & \ddots &  \\
%		0 & a_{d-2} & (0) & & 1
%	\end{smallmatrix}\right) \cdot f = \left(\begin{smallmatrix}
%		0 & & & & (0) \\
%		1/t & 0 & &  & \\
%		a_1 / t & 0 & 0  & & \\
%		\vdots & \vdots & & \ddots &  \\
%		a_{n-2} / t & 0 & (0) & & 0
%	\end{smallmatrix}\right). \]
%	By applying $\kappa$, we get the lemma.
\end{proof}

The restriction map $\g^\vee \twoheadrightarrow (\kfk \ltimes \g_{\geqslant 1})^\vee$ gives rise to a Poisson homomorphism
\[ \g^\vee[e^{-1}] \longrightarrow (\kfk \ltimes \g_{\geqslant 1})^\vee[e^{-1}], \quad \xi \longmapsto \xi|_{\kfk \ltimes \g_{\geqslant 1}}. \]
Consider the fibred product of this restriction map and the quotient map
\[ \g^\vee[e^{-1}] \longrightarrow \g^\vee[e^{-1}] \git (K \ltimes G_{\geqslant 1}), \quad \xi \longmapsto [\xi]_{\bmod K \ltimes G_{\geqslant 1}}, \]
that we denote by
\[ \begin{aligned}
	\delta \colon \g^\vee[e^{-1}] & \longrightarrow (\kfk \ltimes \g_{\geqslant 1})^\vee[e^{-1}] \times \g^\vee[e^{-1}] \git (K \ltimes G_{\geqslant 1}) \\
	\xi & \longmapsto  (\xi|_{\kfk \ltimes \g_{\geqslant 1}}, [\xi]_{\bmod K \ltimes G_{\geqslant 1}})
\end{aligned} \]

\begin{lemma} \label{lemma:gan-ginzburg-decomposition}
	The homomorphism of schemes
	\[ \begin{aligned}
		\psi \colon (K \ltimes G_{\geqslant 1}) \times S_f & \longrightarrow (\kfk \ltimes \g_{\geqslant 1})^\vee[e^{-1}] \times \g^\vee[e^{-1}] \git (K \ltimes G_{\geqslant 1}) \\
		(g,\xi) & \longmapsto  (g \cdot \xi|_{\kfk \ltimes \g_{\geqslant 1}}, [g \cdot \xi]_{\bmod K \ltimes G_{\geqslant 1}})
	\end{aligned} \]
	is an isomorphism.
\end{lemma}

\begin{proof}
	Because the grading is good, we have the inclusions $\kfk \ltimes \g_{\geqslant 1} \subseteq \g_{\geqslant 0} \subseteq [\g, e]$. Then, we get the inclusions $[\g, e]^\perp \subseteq (\g_{\geqslant 0})^{\perp} \subseteq (\kfk \ltimes \g_{\geqslant 1})^{\perp}$.	Therefore, for any $\xi$ in $S_f = - \chi + [\g, e]^\perp$ and for any $g$ in $K \ltimes G_{\geqslant 1}$, we have the equality:
	\[ g \cdot \xi|_{\kfk \ltimes \g_{\geqslant 1}} = - g \cdot \overline{\chi}'. \]
	Therefore, we can simplify the expression for $\psi$:
	\[ \psi(g, \xi) = (- g \cdot \overline{\chi}', [\xi]_{\bmod K \ltimes G_{\geqslant 1}}). \]
	By \Cref{lemma:open-orbit} and \Cref{theorem:slodowy-git}, we see that $\psi$ is an isomorphism because it is the product of the isomorphisms
	\[ K \ltimes G_{\geqslant 1} \cong (\kfk \ltimes \g_{\geqslant 1})^\vee[e^{-1}] \quad \text{and} \quad S_f \cong \g^\vee[e^{-1}] \git (K \ltimes G_{\geqslant 1}). \]
\end{proof}

We can state and prove the decomposition theorem for the minimal Slodowy slice.

\begin{theorem} \label{theorem:decomposition-slodowy}
	The scheme homomorphism
	\[ 	\delta \colon \g^\vee[e^{-1}] \longrightarrow (\kfk \ltimes \g_{\geqslant 1})^\vee[e^{-1}] \times \g^\vee[e^{-1}] \git (K \ltimes G_{\geqslant 1}) \]
	is a Poisson isomorphism, where the Poisson structure on the right-hand side product is the product of Poisson structures.
\end{theorem}

\begin{proof}
	We start by proving that $\delta$ is compatible with the Poisson structures, and then we prove it is an isomorphism.
	
	Both Poisson algebras
	\[  (\Sym(\kfk \ltimes \g_{\geqslant 1}))[e^{-1}] \quad \text{and} \quad \big((\Sym \g)[e^{-1}]\big)^{K \ltimes G_{\geqslant 1}} \]
	are Poisson subalgebras of $(\Sym \g)[e^{-1}]$. They mutually Poisson-commute because of the following equalities of subalgebras in $(\Sym \g)[e^{-1}]$:
	\[ \begin{aligned}
		\big((\Sym \g)[e^{-1}]\big)^{K \ltimes G_{\geqslant 1}} & = \big((\Sym \g)[e^{-1}]\big)^{\kfk \ltimes \g_{\geqslant 1}} \\
		& = \big\{ F \in (\Sym \g)[e^{-1}] \quad | \quad \text{for all} ~  x \in \kfk \ltimes \g_{\geqslant 1},  ~   \{ x, F \} = 0 \big\} \\
		& = \big\{ F \in (\Sym \g)[e^{-1}] \quad | \quad  \{\Sym(\kfk \ltimes \g_{\geqslant 1})[e^{-1}], F \} = 0\big\}. 
	\end{aligned} \]
	\begin{itemize}
		\item The first equality is justified by the connectedness of the group $K \ltimes G_{\geqslant 1}$. 
		\item The second equality is because the Poisson bracket $\{ x, \bullet \}$ coincides with the infinitesimal action of $x$ in $\kfk \ltimes \g_{\geqslant 1}$.
		
		\item The third equality follows because Poisson commuting with $\Sym(\kfk \ltimes \g_{\geqslant 1})[e^{-1}]$ is equivalent to Poisson commuting with $\kfk \ltimes \g_{\geqslant 1}$, since this space generates the non-localised algebra and the Poisson bracket with $e^{-1}$ is given by the formula: $\{e^{-1}, \bullet\} = - e^{-2} \{e, \bullet\}$.
	\end{itemize}
	
	Therefore, the natural algebra homomorphism
	\[ \delta^* \colon \Sym(\kfk \ltimes \g_{\geqslant 1})[e^{-1}] \otimes \big((\Sym \g)[e^{-1}]\big)^{K \ltimes G_{\geqslant 1}} \longrightarrow (\Sym \g)[e^{-1}] \]
	given by tensoring the inclusion homomorphisms is Poisson. So, $\delta$ is a Poisson variety homomorphism. 
	
	To prove it is an isomorphism, we notice that we have the following commutative diagram:
	\[ \begin{tikzcd}[column sep = small]
		& (K \ltimes G_{\geqslant 1}) \times S_f \arrow[ld, "\phi'"', "\sim"'{sloped}] \arrow[rd, "\psi", "\sim"'{sloped}] & \\
		\g^\vee[e^{-1}] \arrow[rr, "\delta"'] & & (\kfk \ltimes \g_{\geqslant 1})^\vee[e^{-1}] \times \g^\vee[e^{-1}] \git (K \ltimes G_{\geqslant 1}).
	\end{tikzcd} \]
	Since both $\phi'$ and $\psi$ are isomorphisms (see \Cref{proposition:dilation-gan-ginzburg} and \Cref{lemma:gan-ginzburg-decomposition}), $\delta$ is an isomorphism by composition.
\end{proof}